\documentclass[12pt]{article}
\usepackage[utf8]{inputenc}
\usepackage[T1]{fontenc}
\usepackage[margin=1in]{geometry}
\usepackage{setspace}
\usepackage{xcolor}
\usepackage[colorlinks=true, urlcolor=blue!60!black, linkcolor=blue!60!black]{hyperref}
\setlength{\parindent}{2em}
\setlength{\parskip}{1.25em}
\renewcommand{\baselinestretch}{1.0}
\begin{document}
\begin{center}
{\LARGE \textit{United States v. Hemani}, No. 24-1234 (2026)\par}
\vspace{0.8em}
{\large Case Brief\par}
\end{center}
\vspace{1em}

\section{Memory Jogger}

A unanimous Supreme Court held that prosecuting a regular marijuana user under 18 U.S.C. \S~922(g)(3) solely for possessing a firearm at home violates the Second Amendment because historical ``habitual drunkard'' laws targeted incapacitated persons for different reasons and only after pre-deprivation process (\href{https://www.courtlistener.com/opinion/10876933/united-states-v-hemani/}{CourtListener}; \href{https://www.supremecourt.gov/opinions/25pdf/24-1234_g2bh.pdf}{Supreme Court Slip Op.}).

\section{Detailed Case Facts}

Ali Danial Hemani is a dual citizen of the United States and Pakistan who was born in Texas, lived most of his life in the Dallas area with his parents, and maintained steady employment (\href{https://supreme.justia.com/cases/federal/us/608/24-1234/}{608 U.S. \_\_\_ (2026), slip op. at 2}). In August 2022, federal agents suspecting Hemani and members of his family of terrorism-related activities executed a search warrant at the family home (\href{https://supreme.justia.com/cases/federal/us/608/24-1234/}{slip op. at 2}; \href{https://ballotpedia.org/United_States_v._Hemani}{Ballotpedia}). Hemani cooperated fully during the search: he surrendered a Glock 19 9mm pistol kept in the house, pointed agents to marijuana on the property, and consented to an interview in which he admitted using marijuana ``about every other day'' (\href{https://supreme.justia.com/cases/federal/us/608/24-1234/}{slip op. at 2}; \href{https://ballotpedia.org/United_States_v._Hemani}{Ballotpedia}). Agents also discovered cocaine in his parents' closet; Hemani claimed ownership of it but stated that his mother had hidden it from him and that he had not used any recently (\href{https://supreme.justia.com/cases/federal/us/608/24-1234/}{slip op. at 2}).

More than six months later, in February 2023, a federal grand jury returned a single-count indictment against Hemani under 18 U.S.C. \S~922(g)(3)---a provision of the Gun Control Act of 1968 that bars any ``unlawful user of or addicted to any controlled substance'' (incorporating all five schedules of the Controlled Substances Act, 21 U.S.C. \S\S~802, 812) from possessing a firearm, punishable by up to 15 years' imprisonment under \S~924(a)(8) and lifetime disarmament under \S~922(g)(1) (\href{https://supreme.justia.com/cases/federal/us/608/24-1234/}{slip op. at 1--3}; \href{https://ballotpedia.org/United_States_v._Hemani}{Ballotpedia}). The government brought no terrorism, cocaine-possession, or drug-trafficking charges; it relied solely on Hemani's admitted marijuana use every other day, without alleging that he was addicted, intoxicated while possessing the gun, or dangerous to himself or others (\href{https://supreme.justia.com/cases/federal/us/608/24-1234/}{slip op. at 2--3}).

\section{Procedural History}

Hemani moved to dismiss the indictment in the United States District Court for the Eastern District of Texas (No. 4:23-cr-00018), arguing that enforcing 18 U.S.C. \S~922(g)(3) against him violated the Second Amendment (\href{https://supreme.justia.com/cases/federal/us/608/24-1234/}{slip op. at 3; Thomas, J., concurring, slip op. at 2}; \href{https://ballotpedia.org/United_States_v._Hemani}{Ballotpedia}). The district court granted the motion and dismissed the indictment (\href{https://supreme.justia.com/cases/federal/us/608/24-1234/}{slip op. at 3}). On appeal, the United States Court of Appeals for the Fifth Circuit affirmed in an unpublished decision (\textit{United States v. Hemani}, No. 24-40137, 2025 WL 354982 (5th Cir. Jan. 2025)), holding that circuit precedent in \textit{United States v. Connelly}, 117 F.4th 269 (5th Cir. 2024), barred applying \S~922(g)(3) to a sober regular marijuana user (\href{https://ballotpedia.org/United_States_v._Hemani}{Ballotpedia}; \href{https://www.supremecourt.gov/DocketPDF/24/24-1234/395080/20260130130853936_Final\%20Hemani\%20NACDL\%20Amicus_1_30_26.pdf}{NACDL Amicus Brief}).

The United States petitioned for a writ of certiorari on June 2, 2025; the Supreme Court granted review on October 20, 2025 (607 U.S. 992), heard oral argument on March 2, 2026 (Principal Deputy Solicitor General Sarah M. Harris for the United States; Erin E. Murphy for Hemani), and issued its decision on June 18, 2026 (reported at 608 U.S. \_\_\_, 146 S. Ct. 1677) (\href{https://supreme.justia.com/cases/federal/us/608/24-1234/}{Justia}; \href{https://www.scotusblog.com/cases/united-states-v-hemani/}{SCOTUSblog}; \href{https://www.cato.org/sites/cato.org/files/2026-09/2025-2026-Supreme-Court-Review-7.pdf}{Cato Supreme Court Review}).

\section{Judicial Votes}

The Supreme Court unanimously affirmed the Fifth Circuit's judgment, \textbf{9--0}, with a seven-Justice majority opinion and three separate concurrences (\href{https://supreme.justia.com/cases/federal/us/608/24-1234/}{Justia}; \href{https://www.scotusblog.com/cases/united-states-v-hemani/}{SCOTUSblog}).

\begin{itemize}
\item \textbf{Majority (7):} Justice Neil M. Gorsuch (author), joined by Chief Justice John G. Roberts, Jr., and Associate Justices Clarence Thomas, Sonia Sotomayor, Brett M. Kavanaugh, Amy Coney Barrett, and Ketanji Brown Jackson (\href{https://supreme.justia.com/cases/federal/us/608/24-1234/}{Justia}).
\item \textbf{Concurrence (1):} Justice Clarence Thomas (joined the majority opinion in full; wrote separately to question \S~922(g)'s validity under the Commerce Clause) (\href{https://supreme.justia.com/cases/federal/us/608/24-1234/}{Justia}).
\item \textbf{Concurrence (2):} Justice Ketanji Brown Jackson, joined by Justice Sonia Sotomayor (both joined the majority opinion in full; wrote separately to criticize \textit{Bruen} and advocate returning to means-end scrutiny) (\href{https://supreme.justia.com/cases/federal/us/608/24-1234/}{Justia}).
\item \textbf{Concurrence in the Judgment (2):} Justice Samuel A. Alito, Jr., joined by Justice Elena Kagan (agreed only with the judgment affirming dismissal on the narrower ground that the government's historical analogues failed as applied to a marijuana user like Hemani) (\href{https://supreme.justia.com/cases/federal/us/608/24-1234/}{Justia}).
\item \textbf{Dissents (0):} None.
\end{itemize}

\section{Holding}

\textbf{Judgment.} Affirmed. The dismissal of the indictment against Ali Hemani was affirmed (\href{https://supreme.justia.com/cases/federal/us/608/24-1234/}{slip op. at 19}).

\textbf{Rule Announced.} Applying \textit{New York State Rifle \& Pistol Ass'n v. Bruen}, 597 U.S. 1 (2022), and \textit{United States v. Rahimi}, 602 U.S. 680 (2024), the government may not automatically strip an individual of his Second Amendment right to possess a firearm under 18 U.S.C. \S~922(g)(3)'s ``unlawful user'' provision solely because he regularly uses a controlled substance such as marijuana, without any pre-deprivation process or proof of incapacitation or dangerousness (\href{https://supreme.justia.com/cases/federal/us/608/24-1234/}{slip op. at 3--7, 18--19}). Historical ``habitual drunkard'' laws (vagrancy, civil-commitment, and surety-of-good-behavior statutes) are not ``relevantly similar'' in either their ``why'' or their ``how,'' because they targeted severely incapacitated persons for different purposes and generally provided judicial process prior to any deprivation of liberty (\href{https://supreme.justia.com/cases/federal/us/608/24-1234/}{slip op. at 7--15}).

\textbf{Holding vs. Dicta.} The narrow holding invalidates only the government's prosecution of Hemani under \S~922(g)(3)'s ``unlawful user'' prong based solely on regular marijuana use without further proof (\href{https://supreme.justia.com/cases/federal/us/608/24-1234/}{slip op. at 18--19}). The Court expressly reserved---and treated as outside its holding---(1) bans on firearm possession by drug ``addicts'' or individuals presently intoxicated; (2) targeted prophylactic laws for specific drugs proven to create a special risk of firearm misuse; (3) felon-in-possession bans under \S~922(g)(1) and mental-commitment bans under \S~922(g)(4); and (4) future \S~922(g)(3) prosecutions supported by individualized proof of dangerousness or proof that a particular drug invariably renders users dangerous (\href{https://supreme.justia.com/cases/federal/us/608/24-1234/}{slip op. at 15 n.6, 18--19}). Statements questioning the government's pretextual dangerousness rationale in light of DOJ non-enforcement memoranda and the 2026 Schedule III rescheduling of some marijuana products reinforce the holding, while Footnote 6's reassurance that \S\S~922(g)(1) and (g)(4) remain undisturbed operates as cautionary dicta (\href{https://supreme.justia.com/cases/federal/us/608/24-1234/}{slip op. at 15--18}).

\section{Analysis of Opinions}

\subsection{Majority Opinion (Justice Gorsuch)}

Writing for seven Justices, Justice Gorsuch began by confirming the governing two-step framework from \textit{Bruen} and \textit{Rahimi}: because the government conceded that \S~922(g)(3) burdens conduct presumptively covered by the Second Amendment's text, the government bore the burden of showing that its prosecution was consistent with the principles underpinning the Nation's historical tradition of firearm regulation, focusing on the ``why'' (purpose) and ``how'' (operation) of proffered analogues (\href{https://supreme.justia.com/cases/federal/us/608/24-1234/}{slip op. at 3--5}). The government relied exclusively on three categories of Founding-era and nineteenth-century ``habitual drunkard'' laws: vagrancy laws, civil-commitment/guardianship statutes, and surety-of-good-behavior statutes (\href{https://supreme.justia.com/cases/federal/us/608/24-1234/}{slip op. at 6}). Gorsuch systematically dismantled the analogy on four cumulative grounds (\href{https://supreme.justia.com/cases/federal/us/608/24-1234/}{slip op. at 7--18}):

\begin{enumerate}
\item \textbf{Different Targets:} In early America's ``culture of copious drinking''---illustrated by the daily alcohol consumption of John Adams, George Washington, Thomas Jefferson, and the 1787 Philadelphia City Tavern farewell dinner---``habitual drunkard'' did not mean a merely regular drinker, but someone so persistently intoxicated as to lose ordinary reasoning faculties and become incapable of managing his affairs (\href{https://supreme.justia.com/cases/federal/us/608/24-1234/}{slip op. at 8--10}). By contrast, the government's reading of \S~922(g)(3) swept in anyone who regularly used any CSA-scheduled substance outside a prescription---from a mild marijuana gummy user to a husband taking his wife's Ambien or a student using a friend's Adderall---without showing incapacitation (\href{https://supreme.justia.com/cases/federal/us/608/24-1234/}{slip op. at 10--11}).
\item \textbf{Different Purposes (``Why''):} Historical vagrancy laws enforced labor norms and suppressed idleness; civil-commitment statutes protected drunkards and their families from financial ruin; and surety-of-good-behavior laws policed moral scandals like blasphemy or eavesdropping---not categorical violent dangerousness, unlike the surety-of-the-peace laws upheld in \textit{Rahimi} (\href{https://supreme.justia.com/cases/federal/us/608/24-1234/}{slip op. at 11--14}).
\item \textbf{Different Operation (``How''):} Each historical analogue required pre-deprivation judicial process (criminal conviction, probate adjudication, or a magistrate's hearing), whereas \S~922(g)(3) automatically divests Second Amendment rights the moment regular unlawful use begins, with no pre-deprivation process (\href{https://supreme.justia.com/cases/federal/us/608/24-1234/}{slip op. at 14--15}).
\item \textbf{Unproven Dangerousness Rationale:} Section 922(g)(3) borrows the CSA's five schedules, which list drugs for public-health or precursor reasons unrelated to violence (\href{https://supreme.justia.com/cases/federal/us/608/24-1234/}{slip op. at 16}). Moreover, DOJ's own non-enforcement policies, legalization across 40 states, and the federal government's 2026 rescheduling of some marijuana products from Schedule I to Schedule III (91 Fed. Reg. 22714) undercut its assertion that millions of marijuana users are categorically dangerous (\href{https://supreme.justia.com/cases/federal/us/608/24-1234/}{slip op. at 16--18}). Quoting then-Judge Barrett's dissent in \textit{Kanter v. Barr}, 919 F.3d 437, 465 (7th Cir. 2019), Gorsuch warned that letting the government designate any group as dangerous on its bare say-so would ``quickly swallow'' the Second Amendment (\href{https://supreme.justia.com/cases/federal/us/608/24-1234/}{slip op. at 18}).
\end{enumerate}

\subsection{Concurring Opinion (Justice Thomas)}

Joining the majority in full, Justice Thomas wrote alone to flag a separate constitutional defect: 18 U.S.C. \S~922(g) appears to exceed Congress's enumerated power under the Commerce Clause (\href{https://supreme.justia.com/cases/federal/us/608/24-1234/}{Thomas, J., concurring, slip op. at 1}). Lower courts have upheld \S~922(g) based on \textit{Scarborough v. United States}, 431 U.S. 563 (1977), requiring only that a firearm crossed state lines at some point in the past (\href{https://supreme.justia.com/cases/federal/us/608/24-1234/}{slip op. at 1--2}). Thomas argued that under both original meaning and \textit{United States v. Lopez}, 514 U.S. 549 (1995), and \textit{United States v. Morrison}, 529 U.S. 598 (2000), mere intrastate gun possession is non-economic activity, \textit{Scarborough} was strictly a statutory decision rather than a constitutional holding, and a bare jurisdictional hook cannot convert the Commerce Clause into a general police power (\href{https://supreme.justia.com/cases/federal/us/608/24-1234/}{slip op. at 2--6}). Noting that 26 years had passed since the Court last granted Commerce Clause relief, Thomas urged lower courts and the Supreme Court to revisit \S~922(g) in an appropriate case (\href{https://supreme.justia.com/cases/federal/us/608/24-1234/}{slip op. at 6}).

\subsection{Concurring Opinion (Justice Jackson, Joined by Justice Sotomayor)}

Justice Jackson, joined by Justice Sotomayor, joined Gorsuch's opinion in full as a faithful application of \textit{Bruen} and \textit{Rahimi}, but wrote separately to urge retiring ``the failed \textit{Bruen} experiment'' in favor of traditional means-end scrutiny (\href{https://supreme.justia.com/cases/federal/us/608/24-1234/}{Jackson, J., concurring, slip op. at 1, 3}). Jackson argued that \textit{Bruen}'s history-and-tradition test is unworkable---producing sharp circuit splits over felon-disarmament cases under \S~922(g)(1)---and blinds courts to modern public-safety stakes (\href{https://supreme.justia.com/cases/federal/us/608/24-1234/}{slip op. at 1--2 \& n.2}). Pointedly, she observed that both the government and the \textit{Hemani} majority were already speaking in the register of means-end tailoring when questioning whether \S~922(g)(3)'s categorical sweep fit its asserted goal of disarming dangerous persons (\href{https://supreme.justia.com/cases/federal/us/608/24-1234/}{slip op. at 2--3}).

\subsection{Opinion Concurring in the Judgment (Justice Alito, Joined by Justice Kagan)}

In a rare cross-ideological pairing, Justice Alito, joined by Justice Kagan, concurred only in the judgment (\href{https://supreme.justia.com/cases/federal/us/608/24-1234/}{Alito, J., concurring in the judgment, slip op. at 1}; \href{https://www.scotusblog.com/2026/06/hemani-proving-the-reason-not-just-naming-it/}{SCOTUSblog}). Alito agreed that \S\S~922(g)(1) and (g)(4) remain presumptively lawful and that the government's habitual-drunkard analogues failed, but would have affirmed on one narrow ground alone: historical laws reached only individuals whose extreme alcohol abuse rendered them persistently incapacitated, whereas the sparse record showed nothing about the quantity, potency, timing, or impairing effect of Hemani's every-other-day marijuana use (\href{https://supreme.justia.com/cases/federal/us/608/24-1234/}{slip op. at 1--2}). Because marijuana today---widely legalized by states and largely tolerated by federal enforcement---resembles alcohol at the Founding, the mismatch was decisive, and the Court needed to say no more (\href{https://supreme.justia.com/cases/federal/us/608/24-1234/}{slip op. at 2--3}).

\subsection{Dissents}

There were no dissenting opinions.

\section{Examples: Future Applications}

\begin{enumerate}
\item \textbf{Same-Side Application (Unconstitutional Under \textit{Hemani}): Prosecution of a Medical-Marijuana Homeowner.} A Colorado resident who uses state-legal medical marijuana twice a week for chronic back pain keeps a locked shotgun at home for self-defense and has no criminal record or history of violence. Federal prosecutors indict her under 18 U.S.C. \S~922(g)(3) relying solely on her regular marijuana use. Under \textit{Hemani}, the prosecution violates the Second Amendment because the government offers no pre-deprivation process and no individualized evidence of incapacitation or dangerousness (\href{https://supreme.justia.com/cases/federal/us/608/24-1234/}{slip op. at 10--11, 18--19}).

\item \textbf{Same-Side Application (Unconstitutional Under \textit{Hemani}): Prosecution of a Non-Violent Off-Label Prescription User.} A graduate student who owns a handgun for home defense regularly takes a friend's unprescribed Schedule II Adderall several nights a week to study, without any evidence of erratic behavior or armed intoxication. Justice Gorsuch explicitly cited this exact scenario to show the overbreadth of the government's theory, making a \S~922(g)(3) prosecution unconstitutional absent individualized proof of dangerousness (\href{https://supreme.justia.com/cases/federal/us/608/24-1234/}{slip op. at 11}).

\item \textbf{Opposite-Side Application (Likely Constitutional): Armed Intoxication in Public During a Violent Confrontation.} Police arrest a chronic drug user who is actively high on marijuana and carrying a loaded handgun during a street-gang dispute, and the government introduces evidence that his drug use impaired his judgment and caused him to threaten bystanders' physical safety. \textit{Hemani} expressly reserved bans on possessing firearms while presently intoxicated and prosecutions backed by individualized proof of dangerousness---the very distinction the Eighth Circuit relied on post-\textit{Hemani} to uphold a \S~922(g)(3) conviction in \textit{United States v. Baxter} (8th Cir. 2026) (\href{https://supreme.justia.com/cases/federal/us/608/24-1234/}{slip op. at 18--19}; \href{https://www.cato.org/sites/cato.org/files/2026-09/2025-2026-Supreme-Court-Review-7.pdf}{Stevenson, Cato Supreme Court Review}).

\item \textbf{Opposite-Side Application (Likely Constitutional): Judicial Commitment of a Severely Incapacitated Opioid Addict.} After an adversarial state court hearing, a probate judge adjudicates an individual as severely addicted to fentanyl and unable to manage his basic affairs, temporarily suspending his firearm rights during court-ordered rehabilitation. Because this measure targets genuine incapacitation (the ``addicted to'' prong reserved in \textit{Hemani}) and provides pre-deprivation judicial process mirroring Founding-era civil-commitment and guardianship laws, it aligns closely with both the ``why'' and ``how'' of historical tradition (\href{https://supreme.justia.com/cases/federal/us/608/24-1234/}{slip op. at 2 n.2, 14--15, 18}).

\item \textbf{Fence-Sitter Application (Close Case): Frequent Methamphetamine User Arrested With a Gun After a Non-Violent Chase.} A defendant who admits using methamphetamine three to four times per week flees police in a vehicle and is arrested with a pistol and drug paraphernalia, but prosecutors introduce no expert or factual evidence proving that his methamphetamine use caused violent behavior or persistent incapacitation. As the Eighth Circuit's post-\textit{Hemani} remand in \textit{United States v. Striplin} (8th Cir. 2026) illustrates, this falls squarely in the gray zone left open by \textit{Hemani}: methamphetamine is far more potent and federally disfavored than marijuana, yet without record proof of either categorical drug-induced dangerousness or individualized danger, prosecutors cannot rely on mere frequency of use alone (\href{https://supreme.justia.com/cases/federal/us/608/24-1234/}{slip op. at 18--19}; \href{https://www.cato.org/sites/cato.org/files/2026-09/2025-2026-Supreme-Court-Review-7.pdf}{Stevenson, Cato Supreme Court Review}).
\end{enumerate}

\section{Critique}

\textbf{Originalist and Textualist Perspectives.} Originalist commentators generally praised \textit{Hemani}'s refusal to let the government strip a core constitutional right on its bare ``say-so.'' Professor Stephanie Barclay argued that \textit{Hemani} demonstrates originalism ``at its best'': history fixes the limited set of permissible reasons (such as incapacitation or proven dangerousness) that can justify disarmament, while an as-applied evidentiary record checks against government pretext (\href{https://www.scotusblog.com/2026/06/hemani-proving-the-reason-not-just-naming-it/}{Barclay, SCOTUSblog}). Yet other textualist and criminal-law scholars faulted the majority for ducking statutory vagueness. Professor Joel S. Johnson argued that \S~922(g)(3)'s undefined phrase ``unlawful user'' is a textbook due-process vagueness defect, and that resolving the case through a fact-bound Second Amendment ruling rather than a narrowing statutory construction left lower courts without portable guidance (\href{https://www.scotusblog.com/2026/06/a-victory-for-the-defendant-in-united-states-v-hemani-but-little-guidance-for-the-lower-courts/}{Johnson, SCOTUSblog}). Meanwhile, Justice Thomas's concurrence highlighted a deeper originalist tension: under a strict reading of the Commerce Clause, federal gun-possession bans like \S~922(g) may be ultra vires altogether (\href{https://supreme.justia.com/cases/federal/us/608/24-1234/}{Thomas, J., concurring, slip op. at 1--6}).

\textbf{Progressive and Doctrinal Perspectives.} Progressive and pragmatic critics echoed Justice Jackson's concurrence, arguing that \textit{Hemani} exposes the incoherence of \textit{Bruen}'s history-and-tradition test (\href{https://supreme.justia.com/cases/federal/us/608/24-1234/}{Jackson, J., concurring, slip op. at 1--3}). Professor Dru Stevenson criticized the majority for effectively engaging in covert overbreadth and means-end scrutiny while ignoring DOJ's January 2026 interim final rule narrowing 27 C.F.R. \S~478.11 and throwing the FBI's National Instant Criminal Background Check System (NICS) into administrative disarray (\href{https://www.cato.org/sites/cato.org/files/2026-09/2025-2026-Supreme-Court-Review-7.pdf}{Stevenson, Cato Supreme Court Review}). Second Amendment scholars affiliated with the Duke Center for Firearms Law similarly noted that \textit{Hemani}'s hyper-narrow holding invites protracted, case-by-case litigation over drug potency and individualized dangerousness, collapsing \textit{Bruen}'s historical inquiry into ad hoc evidentiary battles (\href{https://firearmslaw.duke.edu/2026/03/breaking-down-the-hemani-arguments}{Duke Center for Firearms Law}; \href{https://www.cato.org/sites/cato.org/files/2026-09/2025-2026-Supreme-Court-Review-7.pdf}{Stevenson, Cato Supreme Court Review}).

\textbf{Underlying Competing Values.} \textit{Hemani} reflects a clash between \textit{individual constitutional liberty and notice}---preventing automatic, permanent disarmament of millions of otherwise law-abiding citizens based on shifting federal drug schedules---and \textit{legislative deference and administrability}, which favor bright-line prophylactic rules in firearms background checks and public-safety enforcement (\href{https://www.scotusblog.com/2026/06/hemani-proving-the-reason-not-just-naming-it/}{Barclay, SCOTUSblog}; \href{https://www.cato.org/sites/cato.org/files/2026-09/2025-2026-Supreme-Court-Review-7.pdf}{Stevenson, Cato Supreme Court Review}).

\section{Key Quotations}

\begin{quote}
``And the habitual drunkard laws on which the government relies here differ dramatically from \S~922(g)(3)'s unlawful user provision on every single metric the government invites us to consider: They targeted different kinds of people, did so for different purposes, and operated in different ways.''
\end{quote}
\noindent --- Justice Neil M. Gorsuch, Majority Opinion, slip op. at 7 (\href{https://supreme.justia.com/cases/federal/us/608/24-1234/}{Justia}; \href{https://www.supremecourt.gov/opinions/25pdf/24-1234_g2bh.pdf}{Slip Op.}; \href{https://www.courtlistener.com/opinion/10876933/united-states-v-hemani/}{CourtListener}).

\begin{quote}
``On the government's account, the statute automatically divests an individual of his constitutional right to bear arms the moment he becomes an unlawful user and until he ends his drug use---all without any pre-deprivation process.''
\end{quote}
\noindent --- Justice Neil M. Gorsuch, Majority Opinion, slip op. at 15 (\href{https://supreme.justia.com/cases/federal/us/608/24-1234/}{Justia}; \href{https://www.supremecourt.gov/opinions/25pdf/24-1234_g2bh.pdf}{Slip Op.}).

\begin{quote}
``All of which leaves it awkwardly positioned to suggest that the millions of Americans who now regularly use marijuana are categorically and unusually dangerous.''
\end{quote}
\noindent --- Justice Neil M. Gorsuch, Majority Opinion, slip op. at 17 (\href{https://supreme.justia.com/cases/federal/us/608/24-1234/}{Justia}; \href{https://www.supremecourt.gov/opinions/25pdf/24-1234_g2bh.pdf}{Slip Op.}).

\begin{quote}
``As a matter of both original meaning and this Court's precedents, \S~922(g)(3) appears to exceed Congress's enumerated power to regulate interstate commerce.''
\end{quote}
\noindent --- Justice Clarence Thomas, Concurring Opinion, slip op. at 1 (\href{https://supreme.justia.com/cases/federal/us/608/24-1234/}{Justia}; \href{https://www.supremecourt.gov/opinions/25pdf/24-1234_g2bh.pdf}{Slip Op.}).

\begin{quote}
``In a future case that squarely presents the question, we should consider whether to retire the failed Bruen experiment and return to an explicit assessment of Congress's ends and means when deciding the constitutionality of firearm restrictions.''
\end{quote}
\noindent --- Justice Ketanji Brown Jackson, Concurring Opinion, slip op. at 3 (\href{https://supreme.justia.com/cases/federal/us/608/24-1234/}{Justia}; \href{https://www.supremecourt.gov/opinions/25pdf/24-1234_g2bh.pdf}{Slip Op.}).

\end{document}
